\section{Structural limits}\label{sec:limits}

The algorithms of Sections~\ref{sec:grids}--\ref{sec:constraints} rest on
four hypotheses: a product domain (or a feasible rounding law), an upper
bound $L$ on coordinate curvature, quadratic growth at a unique minimizer,
and a correction that charges for rounding every coordinate. This section
shows that each of them does real work, and that the dependence of
Theorem~\ref{thm:approx} on its parameters is close to necessary.
\begin{itemize}
\item A unique minimizer without a quantitative growth bound does not help:
box QP with bag size three and a unique minimizer is NP-hard on instances
with $\log\kappa=O(I)$, so $\kappa$ cannot be replaced by $\log\kappa$ unless
P${}={}$NP (Proposition~\ref{lim:prop:unique}).
\item In the value-oracle model the exponent $p/2$ of $\kappa$ is optimal
(Proposition~\ref{lim:prop:oracle}), and under the exponential-time
hypothesis the dependence on $p$ is exponential even when $\kappa\le2$, and the exponent of $\kappa$ must grow linearly with $p$
(Propositions~\ref{prop:lbwidth} and~\ref{prop:lbproduct}).
\item Exact piecewise-quadratic messages can need exponentially many pieces
at bag size three and $\kappa\le80$ (Proposition~\ref{lim:prop:messages}).
\item Growth towards a continuum of minimizers does not bound the size of
any corrected-grid certificate (Proposition~\ref{lim:prop:setgrowth}); with
two isolated minimizers, the single-center grids of
Section~\ref{sec:growth} need exponential time
(Proposition~\ref{prop:twocenters}).
\item A path of local affine equalities makes the problem NP-hard at bag size
three and $\kappa=1$ (Proposition~\ref{lim:prop:constraints}).
\item Bag-local corrections are not valid lower bounds
(Proposition~\ref{prop:star}), and local moment consistency of any finite
order does not give an error bound in which large positive curvature is
harmless (Proposition~\ref{lim:prop:moments}).
\end{itemize}
Del Pia and Khajavirad prove that continuous box QP is strongly NP-hard at
treewidth two, even with bounded integer coefficients, and that quartic box
problems on paths are strongly NP-hard \cite{DelPiaKhajavirad2026}. Their
instances need not have a unique minimizer, so their result does not by
itself show that $\kappa$ must enter the running time.

\subsection{Conditioning cannot be dropped}\label{sec:limits-conditioning}

The following construction refines the simple weak NP-hardness reduction
in \cite[Remark~2]{DelPiaKhajavirad2026}: a stronger concave penalty puts
every minimizer at a vertex, a lexicographic linear term makes the minimizer
unique, and the growth constant is computed explicitly.

\begin{proposition}[Unique optimizers without conditioning]\label{lim:prop:unique}
Let $a_1,\ldots,a_n$ be positive integers, $n\ge2$, and let $T$ be an
integer with $0\le T\le U:=\sum_ia_i$. Put $M=\sum_ia_i^2$ and
$\varepsilon=1/(n2^{n+1})$. On the box
\[
x\in[0,1]^n,\qquad s=(s_1,\ldots,s_{n-1})\in[-U,2U]^{n-1},
\qquad s_0:=0,\quad s_n:=T,
\]
minimize
\[
F(x,s)=\sum_{i=1}^n(s_i-s_{i-1}-a_ix_i)^2
+M\sum_{i=1}^nx_i(1-x_i)+\varepsilon\sum_{i=1}^n2^{i-1}x_i .
\]
Then:
\begin{enumerate}
\item The bags $\{s_{i-1},s_i,x_i\}$, with the constants $s_0,s_n$ omitted
  and arranged as a path in the order $i=1,\ldots,n$, form a tree
  decomposition with $p\le3$. All data have binary length polynomial in
  that of $(a,T)$.
\item $F$ has a unique minimizer $z^*=(x^*,s^*)$, and $x^*\in\{0,1\}^n$.
  If some $v\in\{0,1\}^n$ satisfies $a^{\T}v=T$, then
  $a^{\T}x^*=T$ and $\OPT<1/(2n)$; otherwise $\OPT\ge1/n$.
\item $L=4$ is a valid upper coordinate-curvature bound and
  \[
  g=\frac{\varepsilon}{n\,(1+2(n-1)M)}
  \]
  is a valid growth constant. Hence
  $\kappa\le2^{n+3}n^2(1+2nM)$ and $\log_2\kappa=O(I)$.
\end{enumerate}
\end{proposition}

\begin{proof}
Part 1 is immediate: $s_i$ lies in the consecutive bags $i$ and $i+1$, each
$x_i$ in one bag, and each square and unary term lies in a bag.

\emph{Elimination of the states.} For $x\in[0,1]^n$ let
$e(x)=a^{\T}x-T$, $\sigma_0(x)=0$, and
\[
\sigma_k(x)=\sum_{j\le k}a_jx_j-\frac kne(x)\qquad(k=1,\ldots,n),
\]
so that $\sigma_n(x)=T$. For a feasible $s$ put $\delta_k=s_k-\sigma_k(x)$,
with $\delta_0=\delta_n=0$. Then
$s_i-s_{i-1}-a_ix_i=-e(x)/n+\delta_i-\delta_{i-1}$, and
$\sum_i(\delta_i-\delta_{i-1})=0$. Expanding the squares gives the identity
\begin{equation}\label{lim:eq:split}
F(x,s)=\Phi(x)+\sum_{i=1}^n(\delta_i-\delta_{i-1})^2,\qquad
\Phi(x)=\frac{e(x)^2}n+M\sum_ix_i(1-x_i)+\varepsilon\sum_i2^{i-1}x_i .
\end{equation}
Because $0\le\sum_{j\le k}a_jx_j\le U$ and $|e(x)|\le U$, the vector
$(\sigma_1(x),\ldots,\sigma_{n-1}(x))$ lies in $[-U,2U]^{n-1}$.
Since $\delta_k=\sum_{i\le k}(\delta_i-\delta_{i-1})$, the Cauchy--Schwarz
inequality gives $\delta_k^2\le k\sum_i(\delta_i-\delta_{i-1})^2$. Summing
over $k\le n-1$ gives the path inequality
\begin{equation}\label{lim:eq:path}
\sum_{i=1}^n(\delta_i-\delta_{i-1})^2\ge\frac2{n(n-1)}\sum_{k=1}^{n-1}
\delta_k^2\qquad(\delta_0=\delta_n=0).
\end{equation}

\emph{The reduced function.} Its Hessian is
$\nabla^2\Phi=\tfrac2naa^{\T}-2MI\preceq0$, since
$\norm a^2=M$. Thus $\Phi$ is concave. At a vertex $v$,
$\Phi(v)=e(v)^2/n+\varepsilon b(v)$ with $b(v)=\sum_i2^{i-1}v_i$, a
bijection of $\{0,1\}^n$ onto $\{0,\ldots,2^n-1\}$. If $\Phi(v)=\Phi(v')$,
then $(e(v)^2-e(v')^2)/n=\varepsilon(b(v')-b(v))$; the left side lies in
$n^{-1}\Z$ and the right side has absolute value less than
$\varepsilon2^n=1/(2n)$, so both vanish and $v=v'$. Let $x^*$ be the unique
minimizing vertex, $\Phi^*=\Phi(x^*)$, and
$\gamma=\min_{v\ne x^*}(\Phi(v)-\Phi^*)$. We claim $\gamma\ge\varepsilon$.
If $e(v)^2=e(x^*)^2$, the difference is a positive integer multiple of
$\varepsilon$. Otherwise $e(v)^2-e(x^*)^2$ is a nonzero integer. It cannot be
negative, since then $\Phi(v)-\Phi^*\le-1/n+\varepsilon(2^n-1)<0$. Hence it
is at least one and $\Phi(v)-\Phi^*\ge1/n-\varepsilon(2^n-1)>1/(2n)\ge
\varepsilon$.

For $x\in[0,1]^n$ let $Y$ have independent coordinates with
$\Pr(Y_i=1)=x_i$. By concavity and Jensen's inequality,
$\Phi(x)=\Phi(\E Y)\ge\E\Phi(Y)\ge\Phi^*+\gamma\Pr(Y\ne x^*)$. Since $x^*$ is
binary, $\Pr(Y_i\ne x_i^*)=|x_i-x_i^*|\le1$, so
$\Pr(Y\ne x^*)\ge\max_i|x_i-x_i^*|\ge\frac1n\norm{x-x^*}^2$. Therefore
\begin{equation}\label{lim:eq:phigrowth}
\Phi(x)-\Phi^*\ge\frac\varepsilon n\norm{x-x^*}^2\qquad(x\in[0,1]^n).
\end{equation}

\emph{Uniqueness and growth.} By \eqref{lim:eq:split} and
\eqref{lim:eq:phigrowth}, $z^*=(x^*,\sigma(x^*))$ is the unique minimizer
and $\OPT=\Phi^*$. Let $z=(x,s)$ be feasible and $h=x-x^*$. Then
$\sigma_k(x)-\sigma_k(x^*)=\sum_jc_{kj}a_jh_j$ with
$c_{kj}=[j\le k]-k/n\in[-1,1]$, so
$\norm{\sigma(x)-\sigma(x^*)}^2\le(n-1)M\norm h^2$ and
\[
\norm{z-z^*}^2\le\norm h^2+2\norm\delta^2+2\norm{\sigma(x)-\sigma(x^*)}^2
\le(1+2(n-1)M)\norm h^2+2\norm\delta^2 .
\]
By \eqref{lim:eq:split}--\eqref{lim:eq:phigrowth},
$F(z)-\OPT\ge(\varepsilon/n)\norm h^2+\frac2{n(n-1)}\norm\delta^2$. Since
$g(1+2(n-1)M)=\varepsilon/n$ and $g\le\varepsilon/n<1/(n(n-1))$, this gives
$F(z)-\OPT\ge g\norm{z-z^*}^2$. On coordinate lines,
$\partial^2F/\partial s_k^2=4$ and $\partial^2F/\partial x_i^2=2a_i^2-2M\le0$,
so $L=4$ is valid and $\kappa\le4/g\le2^{n+3}n^2(1+2nM)$. Both $n$ and
$\log M$ are $O(I)$.

\emph{Decision.} If $a^{\T}v=T$ for a vertex $v$, then
$\OPT\le\Phi(v)\le\varepsilon(2^n-1)<1/(2n)$. As shown above,
$e(x^*)^2\le e(v)^2=0$. If no vertex solves the instance, every vertex has
$e(v)^2\ge1$, and $\OPT=\Phi(x^*)\ge1/n$.
\end{proof}

\begin{corollary}[The dependence on $\kappa$ is not polylogarithmic]
\label{lim:cor:nopolylog}
Unless $\mathrm P=\mathrm{NP}$, no algorithm computes, for every rational
continuous box QP with a supplied tree decomposition of maximum bag size at
most three and a unique minimizer, the optimal value within $2^{-q}$ in time
polynomial in $I+q+\log\kappa$. In particular, no algorithm achieves the
guarantee of Theorem~\ref{thm:approx} with $f(3,\kappa)$ bounded by a
polynomial in $\log\kappa$.
\end{corollary}

\begin{proof}
Apply such an algorithm to the instances of
Proposition~\ref{lim:prop:unique} with $q=\lceil\log_2(4n)\rceil$. An
approximation within $1/(4n)$ separates $\OPT<1/(2n)$ from $\OPT\ge1/n$,
and $\log\kappa=O(I)$. This would decide Subset Sum, which is NP-complete
\cite{GareyJohnson1979}, in polynomial time.
\end{proof}

Theorem~\ref{thm:approx} bounds $f(p,\kappa)$, for fixed $p$, by a
polynomial in $\kappa$: the state cap is $O(\sqrt\kappa\log(n+2))$ per
coordinate, and the arithmetic is polynomial in its numerator lengths.
Hence for fixed $p$ the unique-optimum problem is polynomial on instances
with $\kappa\le I^{O(1)}$ (Theorem~\ref{thm:exact}), while
Proposition~\ref{lim:prop:unique} makes it NP-hard on instances with
$\kappa\le2^{O(I)}$. The parameter $\kappa$ therefore locates the
complexity boundary at fixed width.


\begin{remark}[Conditioning of the bounded-coefficient reduction]
\label{lim:rem:dk}
The strongly NP-hard instances of \cite[Theorem~3]{DelPiaKhajavirad2026}
have bags of size three and coefficients bounded by absolute constants,
but their minimizers need not be unique. Consider their construction for
the two-item instance $a=(B,B-1)$, $T=B$, $B\ge3$, with objective $\Psi$
(their (20)--(24), all weights one) on the unit box. Here
$\ell=\lceil\log_2(2B)\rceil$, $D=2^\ell\ge2B$, and there are $5\ell+2$
variables. Every term of $\Psi$ is nonnegative. A zero of $\Psi$ encodes a
binary solution, and $(1,0)$ is the only one; all copied and state variables
are then determined by the recurrences. Hence $\Psi$ has a unique minimizer
$v^*$ with value zero. By the qualitative growth lemma for quadratics with a
unique minimizer, some $g>0$ is valid.

Encode instead the choice $(0,1)$ with every copy, digit and running-sum
recurrence exact. Only the final square
$(s_2-w_\ell)^2=((B-1-B)/D)^2$ is nonzero, so $\Psi(y')=D^{-2}$. All
$\ell$ copies of both choice bits change by one, so
$\norm{y'-v^*}^2\ge2\ell$ and every valid growth constant satisfies
$g\le1/(2\ell D^2)$. The largest diagonal entry of $\nabla^2\Psi$ is ten,
attained by the digit-state variables. Therefore
\[
\kappa\ge20\ell D^2\ge80\ell B^2 .
\]
For $B=2^m$ there are $5m+7$ variables and coefficients bounded by five, so
$\kappa$ is exponential in the input length, consistent with
Corollary~\ref{lim:cor:nopolylog}.

Negative curvature is also large relative to growth. Encode the fractional
choice $(1/B,1)$; since $B\cdot B^{-1}+(B-1)\cdot1=B$, every square
residual can be kept at zero with all variables in $[0,1]$, and
$\Psi(y)=q/B$ with $q=1-1/B$. The displacement $d=y-v^*$ annihilates every
affine residual, so $d^{\T}\nabla^2\Psi\,d=-2(q^2+1)$, coming from the
two endpoint penalties. With $R=\norm d^2$, every valid growth constant has
$g\le q/(BR)$, while $\nu\ge2(q^2+1)/R$. Hence
\[
\frac\nu g\ge2B\Bigl(q+\frac1q\Bigr)\ge4B .
\]
This bound on $\nu/g$ also holds for arbitrary positive weights on the
squares and a common positive weight on the two endpoint penalties: both
estimates scale with the penalty weight, and the square weights do not
enter. The reduction therefore does not contradict an algorithm
parameterized by $\nu/g$ either.
\end{remark}

\subsection{The exponent of $\kappa$ in the value-oracle model}
\label{sec:limits-oracle}

The corrected-grid lower bound needs only exact factor values. In a model
where factors are available only through evaluations, the exponent $p/2$ of
$\kappa$ cannot be improved. The construction is the classical hidden-well
argument \cite{NemirovskiYudin1983}, adapted to the present hypotheses.

\begin{proposition}[Value-oracle lower bound]\label{lim:prop:oracle}
Let $p\ge1$, $L>0$, $\kappa\ge8p$, $g=L/\kappa$, and $X=[0,1]^p$. For
$c\in X$ put
\[
F_c(x)=\min\Bigl\{gp,\ \frac L2\norm{x-c}^2\Bigr\}.
\]
Each $F_c$ has upper coordinate curvature $L$, the unique minimizer $c$
with $F_c(c)=0$, and growth $F_c(x)\ge g\norm{x-c}^2$ on $X$; a single
bag of size $p$ is a valid decomposition. Suppose a deterministic algorithm
receives $p$, $L$, $\kappa$ and value access to $F_c$ for an unknown $c$,
and always returns $y\in X$ with $F_c(y)<Lp/\kappa$. Then on some $F_c$ it
makes at least $(\kappa/(8p))^{p/2}-1$ evaluations. A randomized algorithm
that makes at most $Q$ evaluations succeeds, for some $c$, with probability
at most $(Q+1)(8p/\kappa)^{p/2}$.
\end{proposition}

\begin{proof}
On a coordinate line, $\frac L2\norm{x-c}^2-\frac L2x_i^2$ is affine in
$x_i$ and $gp-\frac L2x_i^2$ is concave; their minimum is concave. Thus
$t\mapsto F_c-Lt^2/2$ is concave on every coordinate line. For $x\in X$,
$\norm{x-c}^2\le p$, so $gp\ge g\norm{x-c}^2$; and $L/2\ge g$ because
$\kappa\ge2$. This proves growth, and $c$ is the unique zero.

Moreover $F_c(x)=gp$ unless $\norm{x-c}<r$, where $r^2=2gp/L=2p/\kappa$.
Since $\kappa\ge8p$, $2r\le1$. Let $\mathcal C$ consist of the points of
$X$ whose coordinates lie in $\{0,2r,4r,\ldots\}$. Then
$|\mathcal C|=(\lfloor1/(2r)\rfloor+1)^p\ge(2r)^{-p}=(\kappa/(8p))^{p/2}$,
and the open balls $B(c,r)$, $c\in\mathcal C$, are pairwise disjoint.

Run the deterministic algorithm with every answer equal to $gp$. Let
$x^1,\ldots,x^Q$ be its evaluation points and $y$ its output. Each of these
$Q+1$ points lies in at most one ball. If $Q+1<|\mathcal C|$, some
$c\in\mathcal C$ has a ball containing none of them. For this $F_c$ all
answers are $gp$, so the algorithm makes the same evaluations and returns
$y$, with $F_c(y)=gp=Lp/\kappa$. This contradicts the guarantee, so
$Q\ge|\mathcal C|-1$.

For a randomized algorithm, choose $c$ uniformly in $\mathcal C$ and fix
the random bits. The run on the constant answers has at most $Q$
evaluations and one output. If $F_c$ differs from $gp$ at none of the
evaluation points, the run on $F_c$ is the same and succeeds only if the
output lies in $B(c,r)$. Hence at most $Q+1$ choices of $c$ succeed.
Averaging over the random bits gives success probability at most
$(Q+1)/|\mathcal C|$ for a uniform $c$, hence for some $c$.
\end{proof}

\begin{remark}[Comparison with the algorithm, and the dependence on $p$]
\label{lim:rem:oracle}
With one bag ($n=p$, $N=1$), each stage of the capped algorithm evaluates
$F$ on at most $K^p$ grid points. By Lemma~\ref{lem:states} and the
trial schedule, $K=O(\sqrt\kappa\log(p+2))$ in the successful trial, and
the earlier trials cost no more. Thus the exponent $p/2$ of $\kappa$ is
attained, and Proposition~\ref{lim:prop:oracle} shows that it cannot be
improved in this model; the remaining gap is a factor
$(C\sqrt p\log(p+2))^p$ times the number of stages. The lower bound is also
exponential in $p$ once $\kappa\ge32p$.
\end{remark}

In the explicit model the dependence on $p$ must be exponential, even at
bounded conditioning, and the exponent of $\kappa$ must grow with $p$.

\begin{proposition}[Exponential dependence on the width at bounded conditioning]
\label{prop:lbwidth}
For a graph $(V,E)$ with $V=\{1,\dots,n\}$ let
\[
\Psi(x)=2^n\Bigl(-\sum_ix_i+2\sum_{ij\in E}x_ix_j\Bigr)+\sum_i2^{i-1}x_i
+\frac1{2n}\sum_i(x_i^2-x_i),\qquad x\in[0,1]^n .
\]
Then $\Psi$ has a unique minimizer $x^*$, which is binary,
$\lfloor\Psi^*/2^n\rfloor=-\alpha(V,E)$ (the independence number), and
$\Psi(x)-\Psi^*\ge\frac1{2n}\norm{x-x^*}^2$ on $[0,1]^n$. With
$L_i=L=1/n$ this gives $\kappa\le2$ (and $\bar\kappa\le2$). The same holds on
$\{0,1\}^n$. Every tree decomposition of the graph is a decomposition of
$\Psi$. Consequently, assuming ETH, no algorithm computes $\OPT$ within $1/2$
for all box quadratics with $\kappa\le2$ (continuous, or all-binary) in time
$2^{o(p)}I^{O(1)}$.
\end{proposition}

\begin{proof}
On binary points $\Psi=2^n\Phi+\tau$ with $\Phi=-\sum x_i+2\sum_Ex_ix_j$
integral and $\tau=\sum2^{i-1}x_i\in\{0,\dots,2^n-1\}$ injective. A binary
minimizer therefore minimizes $\Phi$ and is unique. If the support $S$ of a
binary $x$ spans an edge, deleting an endpoint $u\in S$ of such an edge changes
$\Phi$ by $1-2\deg_S(u)\le-1$; hence $\min\Phi=-\alpha$ and
$\lfloor\Psi^*/2^n\rfloor=-\alpha$.

For $x\in[0,1]^n$ write $\Psi=M+\frac1{2n}\sum(x_i^2-x_i)$, where $M$ is
multilinear. Let $V$ have independent coordinates with $\Pr[V_i=1]=x_i$. Then
$M(x)=\E M(V)=\E\Psi(V)$, and since $\Psi(v)\ge\Psi^*+1$ for binary
$v\ne x^*$, $M(x)-\Psi^*\ge\Pr[V\ne x^*]$. With
$\delta_i=|x_i-x_i^*|\in[0,1]$,
$\Pr[V=x^*]=\prod_i(1-\delta_i)\le1-\max_i\delta_i$, so
$M(x)-\Psi^*\ge\max_i\delta_i\ge\frac1n\sum_i\delta_i$. Also
$x_i(1-x_i)\le\delta_i$ because $x_i^*\in\{0,1\}$. Hence
$\Psi(x)-\Psi^*\ge(\frac1n-\frac1{2n})\sum_i\delta_i\ge\frac1{2n}\norm{x-x^*}^2$.
The only squared terms are $\frac1{2n}x_i^2$, so $L=1/n$ and $L/g\le2$.

The instance has $O(n^2)$ coefficients of $O(n)$ bits, so $I=n^{O(1)}$, and
the single bag $V$ gives $p=n$. An algorithm as excluded would compute
$\alpha$ in time $2^{o(n)}$, contradicting ETH
\cite{ImpagliazzoPaturiZane2001}.
\end{proof}

Under the strong exponential-time hypothesis the same reduction, applied with
a supplied tree decomposition of the graph, excludes $(2-\epsilon)^pI^{O(1)}$
by the lower bound of \cite{LokshtanovMarxSaurabh2018} for independent set.

For $L=0$ the corrected grids consist of box endpoints, and our algorithm takes
$O(2^p\poly(I))$ time, so in that case the dependence on $p$ is optimal up to the
base of the exponential.

\begin{proposition}[The exponent of the conditioning must grow with the width]
\label{prop:lbproduct}
Assume rETH. There are no computable $f$ and computable nondecreasing
unbounded $\psi$ such that some algorithm, on every integer box quadratic with
supplied decomposition and a unique minimizer, outputs a feasible point within
$1/2$ of $\OPT$ in time $f(p)\,(\kappa I)^{p/\psi(p)}$. In particular a bound
$f(p)\kappa^{a(p)}I^{C}$ requires $a(p)\ne o(p)$; the bound of
Theorem~\ref{thm:approx} has $a(p)=p/2+O(1)$.
\end{proposition}

The proof, a reduction from multicolored clique with randomized isolation,
is in Appendix~\ref{app:lbproduct}. Together with
Proposition~\ref{prop:lbwidth} and Corollary~\ref{lim:cor:nopolylog}, it shows
that neither parameter can be dropped and that the shape $f(p)\kappa^{\Theta(p)}$
cannot be improved to $f(p)\kappa^{o(p)}$ in general. The constant in the
exponent, the necessity of the factor $p^p$, and whether the same holds for
purely continuous instances or with a deterministic reduction, are open.


\subsection{Exact messages can be exponentially large}
\label{sec:limits-messages}

For forests, exact value functions of box QP are piecewise quadratic with
linearly many pieces \cite{DelPiaKhajavirad2026}. At bag size three, exact
messages can be exponentially large even when $\kappa$ is bounded and the
optimizer is unique. This is why the algorithm uses grids rather than exact
piecewise-quadratic messages.

\begin{proposition}[Exponential messages at bounded conditioning]
\label{lim:prop:messages}
For $m\ge2$ let $S_t\in[0,2^t-1]$ and $z_t\in[0,1]$, $t=1,\ldots,m$, put
$S_0:=0$, $r_t=S_t-2S_{t-1}-z_t$, and
\[
G_m(S,z)=S_m^2+\sum_{t=1}^mr_t^2+\frac18\sum_{t=1}^mz_t(1-z_t).
\]
Then:
\begin{enumerate}
\item the bags $\{S_{t-1},S_t,z_t\}$ ($S_0$ omitted) form a path
  decomposition with $p=3$, and $I=O(m^2)$;
\item $G_m(S,z)\ge\frac18(\norm S^2+\norm z^2)$ on the box, so zero is the
  unique minimizer and $g=1/8$ is valid;
\item $L=10$ is the exact maximum coordinate curvature, so $\kappa\le80$;
  moreover $\nabla^2G_m\succeq-\frac14I$;
\item let $\mu(v)=\min\{G_m(S,z):S_m=v\}$ for $v\in[0,2^m-1]$, the
  min-marginal of $S_m$. Up to the term $v^2$, it is the message sent
  across the separator $\{S_m\}$ to a leaf bag $\{S_m\}$ carrying $S_m^2$.
  If $[0,2^m-1]$ is covered by
  $K$ intervals on each of which $\mu$ coincides with a quadratic
  polynomial, or if $\mu=\min_{k\le K}q_k$ for quadratic polynomials $q_k$,
  then $K\ge2^{m-1}$.
\end{enumerate}
By Theorem~\ref{thm:approx}, the corrected-grid algorithm nevertheless
solves this family to accuracy $2^{-q}$ in time polynomial in $m+q$.
\end{proposition}

\begin{proof}
Part 1: $S_t$ lies in bags $t$ and $t+1$, and the endpoints $2^t-1$ have
$t$ bits.

Part 2. On the box $S_t\ge0$ and $z_t\ge0$, so
$S_{t-1}=(S_t-z_t-r_t)/2\le(S_t+|r_t|)/2$. By induction,
\[
0\le S_t\le2^{t-m}S_m+\sum_{j=t+1}^m2^{t-j}|r_j| .
\]
The vector $(2^{t-m})_{t\le m}$ has squared norm at most $4/3$. The map
$y\mapsto(\sum_{j>t}2^{t-j}y_j)_t$ equals $\sum_{l\ge1}2^{-l}J^l$, where $J$
shifts a vector by one place, so its operator norm is at most one. Hence $\norm S\le(2/\sqrt3)S_m+\norm r$ and
$\norm S^2\le\frac83S_m^2+2\norm r^2$. Also
$0\le z_t=S_t-2S_{t-1}-r_t\le S_t+|r_t|$, so
$\norm z^2\le2\norm S^2+2\norm r^2$. Thus
$\norm S^2+\norm z^2\le3\norm S^2+2\norm r^2\le8S_m^2+8\norm r^2\le8G_m$.

Part 3. The second derivatives are $2+8=10$ for $S_t$, $t<m$; $2+2=4$ for
$S_m$; and $2-\frac14$ for $z_t$. Writing $G_m$ as a sum of squares plus
$-\frac18\sum z_t^2$ plus linear terms gives $\nabla^2G_m+\frac14P_z\succeq0$,
where $P_z$ projects onto the controls.

Part 4. Let $W(v)=\mu(v)-v^2
=\min\{\sum_tr_t^2+\frac18\sum_tz_t(1-z_t):S_m=v\}\ge0$. The minimum is
attained by compactness. Hence $W(v)=0$ exactly when some feasible point has
$S_m=v$, all $r_t=0$ and binary $z$. Then $S_t=2S_{t-1}+z_t$ and
$v=\sum_t2^{m-t}z_t\in\{0,\ldots,2^m-1\}$. Conversely, for an integer $j$ in
this range, its binary digits give $S_t=\lfloor j/2^{m-t}\rfloor\in
[0,2^t-1]$. So $W$ vanishes exactly at the $2^m$ integers of
$[0,2^m-1]$.

In the first representation, on an interval with nonempty interior,
$W$ coincides with a polynomial of degree at most two. It is not
identically zero, since $W>0$ off the integers. So it has at most two zeros,
and the interval contains at most two integers. A degenerate interval
contains at most one. Hence $2K\ge2^m$. In the second representation, each
integer $j$ has some $k$ with $q_k(j)=\mu(j)$. Then $q_k-v^2\ge W\ge0$ on the
interval and vanishes at $j$. Since $q_k-v^2\not\equiv0$, it vanishes at at
most two integers, and again $2K\ge2^m$.
\end{proof}

\subsection{Point growth cannot be replaced by set growth}
\label{sec:limits-setgrowth}

\begin{proposition}[Set growth does not bound grids]\label{lim:prop:setgrowth}
For $n\ge2$ let $F(x)=\sum_{i=1}^{n-1}(x_i-x_{i+1})^2$ on $[0,1]^n$, with
the path decomposition of bags $\{x_i,x_{i+1}\}$. The optimal set is the
diagonal $S=\{t\mathbf1:t\in[0,1]\}$, and
$F(x)\ge\frac2{n(n-1)}\dist(x,S)^2$; for $n=2$, $F=2\dist(x,S)^2$. Every
valid $L$ is at least $2$. For any grids $G_i\ni0,1$ and the corrections
$d_i(v)=Lw_i(v)^2/8$,
\[
\beta=\min_{y\in\prod_iG_i}(F-D)(y)\le-\frac L8W_j^2
\qquad\hbox{for every }j,
\]
where $W_j$ is the largest gap of $G_j$. The filtering rule of
Proposition~\ref{prop:filter} never removes an interval. Hence every
corrected-grid certificate of accuracy $\varepsilon$, including its
filtering history, has a last-stage grid with at least
$1+1/(2\sqrt\varepsilon)$ nodes in every coordinate. Its size is not
polynomial in $\log(1/\varepsilon)$, although $L$ divided by the set-growth
constant is at most $2n(n-1)$.
\end{proposition}

\begin{proof}
For $x\in[0,1]^n$, $\dist(x,S)^2\le\sum_i(x_i-x_1)^2$, and
$(x_i-x_1)^2\le(i-1)F(x)$ by Cauchy--Schwarz, which gives the set-growth
bound. For $n=2$, $\dist(x,S)^2=(x_1-x_2)^2/2$. The coordinate second
derivatives are $2$ or $4$, so $L\ge2$.

Fix $j$, and let $v_j\in G_j$ be an endpoint of a largest gap, so
$w_j(v_j)=W_j$. Choose the remaining coordinates outward along the path. If
$v_k$ has been chosen and $i=k\pm1$ is the next index, let $[\alpha,\alpha']$
be a gap of $G_i$ containing $v_k$, of length $\gamma_i$, and let $v_i$ be
its endpoint nearest to $v_k$. Then $(v_k-v_i)^2\le\gamma_i^2/4$ and
$w_i(v_i)\ge\gamma_i$. Every edge of the path is used exactly once, so
\[
F(v)-D(v)\le\sum_{i\ne j}\frac{\gamma_i^2}4-\frac L8W_j^2
-\sum_{i\ne j}\frac L8\gamma_i^2\le-\frac L8W_j^2,
\]
because $L\ge2$.

Every $t\in[0,1]$ is a coordinate of the optimal point $t\mathbf1$. By
Proposition~\ref{prop:cellwise}, every interval $[a,b']$ of every grid satisfies
$\min\{m_i(a),m_i(b')\}\le F(t\mathbf1)=0\le U$ for $t\in[a,b']$, so the
filtering rule keeps it. Thus every stage covers $[0,1]^n$, and the final
gap is at least $U-\beta\ge LW_j^2/8\ge W_j^2/4$ since $U\ge\OPT=0$. A gap
at most $\varepsilon$ forces $W_j\le2\sqrt\varepsilon$, hence
$|G_j|\ge1+1/(2\sqrt\varepsilon)$.
\end{proof}

For finitely many optimal values per coordinate, unions of retained cells
restore a bound (Theorem~\ref{thm:cells}). The proposition concerns a continuum of
optimal coordinate values. It is a limit of corrected-grid certificates, not
of the problem: here $F$ is a sum of squares.

\begin{proposition}[Value-oracle barrier for unknown set growth]\label{prop:oraclebarrier}
Fix $k\ge1$, $X=[0,1]^k$ and the coordinate-curvature bound $L=1$. Consider a
deterministic algorithm that queries exact values, gradients and Hessians of
an unknown $C^2$ function $f$ at adaptively chosen points and outputs either a
rational interval of width at most $\varepsilon$ containing $\min_Xf$, or a
point $\hat x$ with $f(\hat x)-\min_Xf\le\varepsilon$. Suppose it is correct
for every $C^2$ function with $\partial_{ii}f\le1$ on $X$ for all $i$ that
satisfies \eqref{eq:setgrowth} for some $g>0$. Then on $f\equiv0$, for which
$\kappa=1$, it makes at least $\bigl(1/(2\sqrt{6\varepsilon})-1\bigr)^k-1$
queries.
\end{proposition}

\begin{proof}
For $c\in X$ and $0<w\le1$ let
$f_{c,w}(x)=-\frac{w^2}6\bigl(1-|x-c|^2/w^2\bigr)^3$ when $|x-c|\le w$ and
$f_{c,w}(x)=0$ otherwise. In the variable $z=(x-c)/w$ with $\rho=|z|^2$,
\[
\partial_if_{c,w}=w\,z_i(1-\rho)^2,\qquad
\partial_{ij}f_{c,w}=\delta_{ij}(1-\rho)^2-4z_iz_j(1-\rho).
\]
Value, gradient and Hessian vanish on $\rho=1$, so $f_{c,w}$ is $C^2$, and
$\partial_{ii}f_{c,w}=(1-\rho)(1-\rho-4z_i^2)\le1$. The minimizer $c$ is
unique with value $-w^2/6$. Inside the ball,
$f_{c,w}+\frac{w^2}6=\frac{w^2}6\bigl(1-(1-\rho)^3\bigr)\ge\frac{w^2}6\rho
=\frac{|x-c|^2}6$, because $1-(1-\rho)^3-\rho=\rho(1-\rho)(2-\rho)\ge0$.
Outside it the gap is $w^2/6\ge\frac{w^2}{6k}|x-c|^2$. Hence
\eqref{eq:setgrowth} holds with $g=w^2/(6k)$. The zero function satisfies it
with $g=1$ and $S=X$.

Run the algorithm on $f\equiv0$; let $q$ be its number of queries and add
the output point, if any, to the query points. Suppose
$q<\bigl(1/(2\sqrt{6\varepsilon})-1\bigr)^k-1$, and let
$M=\lfloor(q+1)^{1/k}\rfloor+1$. Then $M^k>q+1$ and $M<1/(2\sqrt{6\varepsilon})$.
Divide $X$ into $M^k$ closed subcubes of side $1/M$. The interior of one of
them contains none of the at most $q+1$ points. Let $c$ be its center and
choose $w$ with $\sqrt{6\varepsilon}<w<1/(2M)$. The closed ball of radius $w$
lies in the open subcube, so every oracle answer for $f_{c,w}$ at the
recorded points equals the answer for $0$. By determinism the algorithm
produces the same transcript and output on $f_{c,w}$. An interval of width at
most $\varepsilon$ containing $0$ does not contain $-w^2/6<-\varepsilon$. An
output point outside the ball has gap $w^2/6>\varepsilon$ for $f_{c,w}$.
Either way the algorithm is wrong on an admissible function.
\end{proof}

The filtered-grid algorithm of Section~\ref{sec:growth} uses only exact values
and the curvature bound, yet needs no accuracy-dependent number of states under
point growth. Proposition~\ref{prop:oraclebarrier} shows that no algorithm in
that information model can do the same under growth toward an unknown set,
even with $k=p=1$ and $\kappa=1$. The barrier concerns information, not the
input class: the alternatives are not quadratics. It is the classical
hidden-bump argument of information-based complexity, adapted to set growth.
An algorithm for explicit quadratics must therefore exploit the coefficients,
for example through a certificate such as Lemma~\ref{lem:diagcert}, or the
structure of Section~\ref{sec:endpointset}.

\subsection{Product domains are essential}\label{sec:limits-constraints}

\begin{proposition}[Local affine constraints]\label{lim:prop:constraints}
Let $a_1,\ldots,a_n$ and $B$ be positive integers with $a_i\le B$, and put
$a_0=B$. Consider binary $x_0,\ldots,x_n$, continuous
$y_0,\ldots,y_{n+1}\in[0,1]$, and the constraints
\[
y_0=0,\qquad y_{n+1}=1,\qquad y_{i+1}=y_i+\frac{a_i}Bx_i\quad(i=0,\ldots,n),
\]
with objective $F=2^{n+1}x_0+\sum_{i=1}^n2^{i-1}x_i$. Each constraint lies in
the bag $\{y_i,x_i,y_{i+1}\}$, and these bags form a path with $p=3$. The
feasible set is nonempty, the minimizer $z^*$ is unique, and
$F(z)-\OPT\ge\norm{z-z^*}^2/(2n+3)$ for every feasible $z$. Since $F$ is
affine, every $L>0$ is valid; with $L=1/(2n+3)$, $\kappa=1$. Moreover
$x_0^*=0$ if and only if some subset of $\{a_1,\ldots,a_n\}$ sums to $B$.
Consequently, unless $\mathrm P=\mathrm{NP}$, no algorithm solves such
problems in time $f(p,\kappa)\poly(I)$ for any function $f$, even when the
optimal value is required only within $1/2$.
\end{proposition}

\begin{proof}
Summing the recurrences gives $\sum_{i=0}^na_ix_i=B$. The choice $x_0=1$,
$x_i=0$ ($i\ge1$) is feasible. Since all $a_i>0$, it is the only feasible
choice with $x_0=1$. A choice with $x_0=0$ is feasible exactly when its
items sum to $B$: its partial sums then lie in $[0,1]$. For each binary $x$
the constraints determine $y$. Distinct feasible points therefore have
distinct binary vectors and distinct integer objective values. The dummy
value $2^{n+1}$ exceeds every value $\le2^n-1$ of a genuine subset. Hence
the minimizer is unique, and $x_0^*=0$ exactly for yes-instances. All
$2n+3$ coordinates lie in $[0,1]$, so $\norm{z-z^*}^2\le2n+3$, while
$F(z)-\OPT\ge1$ for $z\ne z^*$. An approximation within $1/2$ separates
$\OPT\le2^n-1$ from $\OPT=2^{n+1}$. The data have polynomial binary length,
and Subset Sum is NP-complete \cite{GareyJohnson1979}.
\end{proof}

In the product setting an affine objective is solved exactly by endpoint
dynamic programming. The difficulty here comes entirely from the
constraints. Independent coordinate rounding, the basis of
Proposition~\ref{prop:cellwise}, does not preserve the running-sum equations.


After scaling $y$ by $B$, the constraint matrix of the continuous variables in
Proposition~\ref{lim:prop:constraints} is the incidence matrix of a path,
which is totally unimodular, and the scaled boxes have width $B$. The
pseudopolynomial first level $K_0=W+1$ in Theorem~\ref{thm:tu-approx} therefore
cannot be removed in general unless P${}={}$NP.

\subsection{Local corrections and local moments}\label{sec:limits-local}
Proposition~\ref{prop:star} shows that a correction charged only to the
coordinates of one bag, while the other coordinates are optimized on the
grid, cannot be valid for any fixed allowance: the missing error grows with
the number of outside coordinates and depends on the separator value. The
valid local alternative is exact or certified conditional recourse
(Theorem~\ref{thm:cr-filter}).

A different route to a local error bound keeps exact positive curvature and
discretizes only the negative part. For a continuous-box quadratic
$F(x)=\frac12x^\T Hx+b^\T x+c$ with $H\succeq-\nu I$, one would like a bound
in which $L$ is replaced by $\nu$. Exact positive energy suffices when all
coordinates are conditioned on one global cell.

\begin{lemma}[Globally conditioned moments]\label{lim:lem:mixture}
Let $(m,Y)$ satisfy $m\in C=\prod_i[l_i',u_i']\subseteq X$,
$\Sigma=Y-mm^\T\succeq0$ and $Y_{ii}\le(l_i'+u_i')m_i-l_i'u_i'$. Then
\[
\tfrac12\langle H,Y\rangle+b^\T m+c\ge\OPT-\frac\nu8\sum_i(u_i'-l_i')^2 .
\]
The same bound, averaged with the corresponding weights, holds for any
convex combination of such pairs indexed by global cells.
\end{lemma}

\begin{proof}
The left side equals
$F(m)+\frac12\langle H,\Sigma\rangle\ge F(m)-\frac\nu2\operatorname{tr}\Sigma$.
Also $\Sigma_{ii}=Y_{ii}-m_i^2\le(u_i'-m_i)(m_i-l_i')\le(u_i'-l_i')^2/4$, and
$F(m)\ge\OPT$. Averaging preserves the inequality.
\end{proof}

A global cell index can take $K^n$ values. Sparse relaxations keep exact
information only on bags and match moments across separators.

\begin{definition}[Local moment relaxation of order $k$]\label{lim:def:moments}
Fix a tree decomposition and a finite partition of every $X_i$ into intervals
(cells). A feasible point consists of probability measures $\mu_t$ on
$X_{B_t}=\prod_{i\in B_t}X_i$, one per bag, such that for every tree edge
$\{t,u\}$ with separator $S$, every product $C$ of separator cells, and every
monomial $x_S^\alpha$ with $\abs\alpha\le k$,
$\int_{\{x_S\in C\}}x_S^\alpha\,d\mu_t=\int_{\{x_S\in C\}}x_S^\alpha\,d\mu_u$.
Its value is $\sum_t\int q_t\,d\mu_t$, where $q_t$ is the sum of factors
assigned to $t$. Point masses at points of $X$ are feasible, so the optimal
value $R_k$ satisfies $R_k\le\OPT$. The occupied width $w_i$ is the largest
length of a cell of coordinate $i$ that carries positive mass.
\end{definition}

\begin{proposition}[Every finite order fails]\label{lim:prop:moments}
For each integer $r\ge1$ and rational $0<h\le1/(4r)$ there is a continuous box
QP $F_{r,h}$ with a tree decomposition of maximum bag size three, a unique
minimizer and $\OPT=1/4$, such that $\nu\le2$ and
$g_r=1/(8r(1+8(2r-1)^2))$ is a valid growth constant, both independent of $h$,
and for every choice of cells in which $[0,rh]$ lies in one cell of width at
most $2rh$ for each continuous coordinate and $0$ and $1$ lie in cells of width
at most $\omega$ for each control,
\[
R_{2r-1}\le\frac14-\frac{2r+1}{16r},\qquad
\sum_iw_i^2\le(4r-2)\max\{2rh,\omega\}^2 .
\]
Consequently there is no function $C$ such that
$\OPT-R_k\le C(p,\bar\nu/g,k)\,\bar\nu\sum_iw_i^2$ for all instances, all
valid bounds $\bar\nu\ge\nu$ and growth constants $g$, all orders $k$, and all
cells.
\end{proposition}

The construction and proof are in Appendix~\ref{app:moments}. The two bag
chains realize the even and the odd part of a binomial law on the separator,
whose moments agree up to order $2r-1$. The example does not contradict an
error bound in terms of $L$, which is of order $r/h^2$ here, and it does not
survive exact matching of separator laws, which glues the bag measures into a
global probability measure. It shows that a $\nu$-based local error bound must
use information beyond finite-order separator moments, for example
inequalities that couple both sides of a separator.
